\section{Gravity-decoupling limits}
\label{sec:decoupling}

\subsection{Dynamical vectors and local scales}
\label{subsec:dynamic-background}

At finite compact volume, every harmonic two-form has finite norm and
supplies a dynamical vector in the reduction of $C_3$
\cite{Cadavid:1995bk}. At the contraction, the graviphoton direction
pairs trivially with the exceptional curves, so the two nonzero rows
of \eqref{eq:x9-effective-vectors} belong to vector multiplets.
A gravity-decoupling limit can instead turn these vectors into
background flavor sources. Whether their constraint survives is
therefore a question about the full kinetic form, not the charge
matrix alone.

Work at fixed eleven-dimensional Planck length. For a divisor
combination $P$ the inverse coupling, up to one common normalization,
is the Hodge norm
\begin{equation}
 \mathcal N_P(J)=\int_{\wideX_9}\omega_P\wedge *_J\omega_P
 =-JP^2+\frac{(J^2P)^2}{4\mathcal V},\qquad
 \mathcal V=\frac{J^3}{6}.
 \label{eq:norm-formula}
\end{equation}
This follows from Lefschetz decomposition of harmonic $(1,1)$-forms
\cite{Candelas:1990pi,Strominger:1990pd}. A retained vector needs a
finite nonzero norm in the local units. A mode of divergent norm has
vanishing canonically normalized coupling to a fixed local current.
Once the decoupled vector is omitted, its moment-map equation is not
imposed on the limiting local theory. All vector mixing must be
included in this test.

We consider K\"ahler paths with
\begin{equation}
 \mathcal V\longrightarrow\infty,\qquad
 J C_1,\ J C_2,\ J|_E\ \hbox{bounded}.
 \label{eq:gravity-decoupling}
\end{equation}
The first condition decouples five-dimensional gravity; the others
keep the local particle and surface scales bounded. These hypotheses
are essential. The finite-volume circle limit of
Section~\ref{subsec:x9-circle-limit} is a different operation.

\subsection{The allowed nef face and retained vectors}
\label{subsec:x9-rigid-limits}

The linear space annihilating the local scales is
\begin{equation}
 H=\{w:wC_1=wC_2=0,\ w|_E=0\}
   =\operatorname{span}(D_0,D_1).
 \label{eq:local-scale-plane}
\end{equation}
It is also the pullback of $H^2(X_9;\QQ)$ to the resolution.

\begin{proposition}\label{prop:local-scale-face}
The intersection of the nef cone with $H$ is
$\mathbb R_{\geq0}D_0+\mathbb R_{\geq0}D_1$.
\end{proposition}
\begin{proof}
Both generators are restrictions of nef toric divisors. The effective
curves $D_1D_8|_{\wideX_9}$ and $D_0D_2|_{\wideX_9}$ have
$(D_0,D_1)$ pairings $(0,1)$ and $(4,0)$, respectively. Hence a
nef combination $\alpha D_0+\beta D_1$ must have
$\alpha,\beta\geq0$, proving the reverse inclusion.
\end{proof}

\begin{lemma}\label{lem:dynamical-subspace}
For a path satisfying \eqref{eq:gravity-decoupling}, each normalized
accumulation direction is a nonzero element
$w=\alpha D_0+\beta D_1$ of this cone and has $w^3>0$.
If $\mathcal N_P$ stays bounded along a subsequence with direction
$w$, then $wP=0$ in $H^4(\wideX_9;\mathbb R)$. Conversely, along
$J=J_*+\Lambda w$, every $P\in\operatorname{ann}(w)$ has
$\mathcal N_P\to-J_*P^2$.
\end{lemma}
\begin{proof}
The local pairings define a linear map with kernel $H$.
Their boundedness bounds the transverse component of $J$, while
divergent volume forces its $H$ component to diverge. Its normalized
limits lie in the nef face. The cubic
\eqref{eq:x9-inherited-cubic} is positive on every nonzero ray there.
Writing $J=s(w+o(1))$, $s\to\infty$, gives
\begin{equation}
 \mathcal N_P=s\left[-wP^2+\frac{3(w^2P)^2}{2w^3}\right]+o(s).
\end{equation}
The Hodge index inequality, extended to nef $w$ by continuity,
is $(w^2P)^2\geq w^3(wP^2)$. The bracket is thus at least
$(w^2P)^2/(2w^3)\geq0$. Boundedness forces $w^2P=wP^2=0$.
The form $wPQ$ is negative semidefinite on $w^2P=0$;
a vector of zero square lies in its radical. Since $w^3>0$,
orthogonality also to $w$ gives $wP=0$ in $H^4$.
For the converse, both $JP^2$ and $J^2P$ are independent of
$\Lambda$ when $wP=0$, and the denominator in
\eqref{eq:norm-formula} diverges.
\end{proof}

\begin{theorem}\label{thm:no-rigid-limit}
No gravity-decoupling path satisfying \eqref{eq:gravity-decoupling}
retains the $E_2$ Cartan gauging of $X_9$. In particular the norms
of $V_1=-D_6$ and $V_2=D_5$ diverge. The maximal retained subspaces
for each leading direction, realized on straight K\"ahler rays,
have the following charge rows on $(u_1,u_2,\beta)$:
\begin{equation}
\begin{array}{c|c|c|c}
 w&\operatorname{ann}(w)&\hbox{nonzero charge row}&\rank Q_{\rm dyn}\\\hline
 D_1&\operatorname{span}(D_2,E)&(1,-1,0)&1\\
 D_0&\operatorname{span}(D_8,E)&\hbox{none}&0\\
 \alpha D_0+\beta D_1,\ \alpha\beta>0&\operatorname{span}(E)&\hbox{none}&0 .
\end{array}
\label{eq:x9-limit-table}
\end{equation}
\end{theorem}
\begin{proof}
The intersection form gives the annihilators in
\eqref{eq:x9-limit-table}; the charge rows follow from
\eqref{eq:x9-full-matrix}. None pairs nontrivially with $R$.
The lemma bounds the retained subspace of any path by the appropriate
annihilator. Further divergent subleading scales may remove vectors,
but cannot restore a root-charged one.
For a straight ray take $J_*=(6,18,16,-11,-8,7)$ in the basis
$\mathcal D$. The limiting norm matrices are
$\operatorname{diag}(14,11)$ on $(D_2,E)$ for $w=D_1$, and
$\left(\begin{smallmatrix}5&-1\\-1&11\end{smallmatrix}\right)$
on $(D_8,E)$ for $w=D_0$; on the interior ray the $E$ norm is
$11$. They are positive definite and realize the listed subspaces.
Neither $D_5$ nor $D_6$ belongs to an annihilator, so their norms
cannot stay bounded on any diverging sequence satisfying the hypotheses.
\end{proof}

Thus a local limit may retain $u_1=u_2$, but never
$u_1=u_2=-\beta$. Other compactifications can retain common
compact surfaces and shared gauging after gravity decouples
\cite{Apruzzi:2019vpe,Cvetic:2023plv}; the theorem is specific to
the cone and local-scale conditions of $X_9$.

\subsection{Current exchange with vector mixing}
\label{subsec:x9-current-exchange}

The inverse kinetic matrix makes the loss of the interaction explicit.
For $J=t^ID_I$ set $L_{IJ}=\kappa_{IJK}t^K$ and $v=Lt$.
Polarization of \eqref{eq:norm-formula} gives
\begin{equation}
 G=-L+\frac{vv^{\mathsf T}}{4\mathcal V},\qquad
 G^{-1}=-L^{-1}+\frac{tt^{\mathsf T}}{2\mathcal V}.
 \label{eq:x9-full-inverse}
\end{equation}
Multiplication verifies the second identity using
$t^{\mathsf T}Lt=6\mathcal V$. Let $C$ be the first two columns
of \eqref{eq:x9-full-matrix}. The exchange coefficient of the
two corresponding currents is $\mathsf H=C^{\mathsf T}G^{-1}C$,
including all six vector fields. On $J=J_*+\Lambda w$,
\begin{equation}
 (1,1)\mathsf H_w(1,1)^{\mathsf T}
 =\begin{cases}
  2/(3\Lambda)+O(\Lambda^{-2}),&w=D_0,\\
  1/\Lambda+O(\Lambda^{-2}),&w=D_1,\\
  2/(5\Lambda)+O(\Lambda^{-2}),&w=D_0+D_1.
 \end{cases}
 \label{eq:x9-root-exchange}
\end{equation}
The full matrices are printed in Appendix~\ref{app:integral-charges}.
For $w=D_1$ the difference current retains a finite coupling,
but the shared root current still decouples. Removing the graviphoton
subtracts $tt^{\mathsf T}/(3\mathcal V)$ from $G^{-1}$; its
contribution to $\mathsf H$ is $O(\Lambda^{-3})$, leaving
\eqref{eq:x9-root-exchange} unchanged.

At every finite $\Lambda$, $\mathsf H$ is positive definite, so
$(Q^{\rm eff}b)^\dagger\mathsf H(Q^{\rm eff}b)$ has kernel
$\CC(-1,-1,1)$. Its limiting matrix has the larger kernel
$\{u_1=u_2\}$ on the $D_1$ ray and $\CC^3$ on the other two.
The vacuum constraint of the limiting local theory therefore need
not be the limit of the finite-coupling constraint. This quadratic
expression only diagnoses the rigid moment maps: the chosen $J_*$
leaves $JC_1=JC_2=1$, so these rays are not themselves Higgs vacua.
At finite Planck mass the full conditions of
\eqref{eq:compact-quaternionic-vacua} are still required.
